\documentclass[12pt,a4paper]{article}
\usepackage[a4paper,left=2.5cm,right=2cm,top=2cm,bottom=2cm]{geometry}
\usepackage{fontspec}
\setmainfont{DejaVu Serif}
\setmonofont{DejaVu Sans Mono}
\usepackage{setspace}
\onehalfspacing
\usepackage{enumitem}
\usepackage{booktabs}
\usepackage{amsmath}
\usepackage{tabularx}
\usepackage{array}
\usepackage{xcolor}
\usepackage{listings}
\usepackage{hyperref}
\hypersetup{hidelinks}
\usepackage{titlesec}
\titleformat{\section}{\normalfont\Large\bfseries}{\thesection.}{0.5em}{}
\titleformat{\subsection}{\normalfont\large\bfseries}{\thesubsection.}{0.5em}{}
\setlength{\parindent}{1.25cm}
\setlength{\parskip}{0pt}
\renewcommand{\contentsname}{Contents}

\definecolor{codebg}{RGB}{246,246,246}
\definecolor{codeframe}{RGB}{215,215,215}
\lstset{
  basicstyle=\ttfamily\small,
  backgroundcolor=\color{codebg},
  frame=single,
  rulecolor=\color{codeframe},
  breaklines=true,
  columns=fullflexible,
  keepspaces=true,
  showstringspaces=false,
  aboveskip=0.8em,
  belowskip=0.8em
}

\begin{document}

\begin{titlepage}
\begin{center}
{\large Ministerul Educației și Cercetării al Republicii Moldova}\\[0.3cm]
{\large Universitatea Tehnică a Moldovei}\\[0.3cm]
{\large Facultatea Calculatoare, Informatică și Microelectronică}\\[0.3cm]
{\large Program: Software Engineering}\\[2.3cm]

{\Large \textbf{Operating Systems}}\\[0.8cm]
{\LARGE \textbf{Laboratory Work No. 6}}\\[0.5cm]
{\Large \textbf{Inside the xv6 File System}}\\[0.35cm]
{\large Track B: Build Your Own Small OS (xv6)}

\vfill

\begin{flushright}
\begin{tabular}{ll}
\textbf{Elaborated by:} & \\
st. gr. FAF-241 & \\
\textbf{Ermicev Sofia} & \\
\\
\textbf{Verified by:} & \\
\textbf{Patricia Reițman} &
\end{tabular}
\end{flushright}

\vfill
{\large Chișinău -- 2026}
\end{center}
\end{titlepage}

\tableofcontents
\newpage

\section{Introduction}
The purpose of this laboratory work is to examine the internal organization of the xv6 file system and connect the concepts of directories, inodes, hard links, block addressing, and file-system limits with the actual xv6 source code. The work also includes two user programs: \texttt{mystat}, which displays inode metadata, and \texttt{find}, which recursively traverses the directory tree.

All experiments were performed in the existing \texttt{\~/xv6-riscv} tree. New user programs were added to the \texttt{user/} directory and registered in the \texttt{UPROGS} list in the Makefile.

\section{Part 1. Reading What xv6 Shows}

The file types were first decoded from \texttt{kernel/stat.h}:
\begin{lstlisting}
$ grep -n "T_DIR\|T_FILE\|T_DEVICE" kernel/stat.h
1:#define T_DIR    1 // Directory
2:#define T_FILE   2 // File
3:#define T_DEVICE 3 // Device
\end{lstlisting}

The root directory was listed in xv6. The beginning of the output was:
\begin{lstlisting}
$ ls
.              1 1 1024
..             1 1 1024
README         2 2 2441
cat            2 3 35376
echo           2 4 34312
...
console        3 25 0
\end{lstlisting}

Each line printed by xv6 \texttt{ls} contains the name, type, inode number, and size in bytes.

A new directory \texttt{a} was then inspected:
\begin{lstlisting}
$ ls a
.              1 26 32
..             1 1 1024
\end{lstlisting}

The directory-entry structure and the maximum stored name length were inspected in \texttt{kernel/fs.h}:
\begin{lstlisting}
struct dirent {
  ushort inum;
  char name[DIRSIZ] __attribute__((nonstring));
};

#define DIRSIZ 14
\end{lstlisting}

\subsection{Observations}
\begin{enumerate}[label=\arabic*.]
\item Type \texttt{1} means directory (\texttt{T\_DIR}), type \texttt{2} means regular file (\texttt{T\_FILE}), and type \texttt{3} means device (\texttt{T\_DEVICE}).
\item The root directory has inode number \textbf{1}. In the root directory, both \texttt{.} and \texttt{..} have inode 1 because \texttt{.} refers to the current directory and the root has no parent above itself, so its \texttt{..} also points back to root.
\item An empty directory \texttt{a} has size 32 bytes and contains exactly two entries: \texttt{.} and \texttt{..}. A \texttt{struct dirent} contains a 2-byte \texttt{ushort inum} and a 14-byte name, so one directory entry occupies \textbf{16 bytes}. Two entries therefore occupy \textbf{32 bytes}.
\end{enumerate}

\section{Part 2. The \texttt{mystat} Program}

The program \texttt{user/mystat.c} was implemented using the xv6 \texttt{stat()} call and \texttt{struct stat}. It converts numeric types to readable names and prints the inode number, link count, and size.

\begin{lstlisting}[language=C]
#include "kernel/types.h"
#include "kernel/stat.h"
#include "user/user.h"

char *
typename(int t)
{
    if(t == T_DIR)
        return "directory";
    if(t == T_FILE)
        return "file";
    if(t == T_DEVICE)
        return "device";
    return "?";
}

int
main(int argc, char *argv[])
{
    struct stat st;

    for(int i = 1; i < argc; i++){
        if(stat(argv[i], &st) < 0){
            fprintf(2, "mystat: cannot stat %s\n", argv[i]);
            continue;
        }

        printf("%s: %s, inode %d, links %d, size %ld bytes\n",
               argv[i], typename(st.type), st.ino, st.nlink,
               (long)st.size);
    }

    exit(0);
}
\end{lstlisting}

The program was added to the Makefile and successfully built. It was first tested on a file, a directory, and a device:
\begin{lstlisting}
$ mystat README . console
README: file, inode 2, links 1, size 2441 bytes
.: directory, inode 1, links 1, size 1024 bytes
console: device, inode 26, links 1, size 0 bytes
\end{lstlisting}

A hard link was then created:
\begin{lstlisting}
$ ln README r2
$ mystat README r2
README: file, inode 2, links 2, size 2441 bytes
r2: file, inode 2, links 2, size 2441 bytes
\end{lstlisting}

After removing the original name, the second name continued to access the same inode:
\begin{lstlisting}
$ rm README
$ mystat r2
r2: file, inode 2, links 1, size 2441 bytes
\end{lstlisting}

The unlink implementation was inspected in \texttt{kernel/sysfile.c}. The relevant part is:
\begin{lstlisting}[language=C]
if (ip->type == T_DIR) {
    dp->nlink--;
    iupdate(dp);
}
...
ip->nlink--;
iupdate(ip);
iunlockput(ip);
\end{lstlisting}

The on-disk inode structure was inspected in \texttt{kernel/fs.h}:
\begin{lstlisting}[language=C]
struct dinode {
  short type;
  short major;
  short minor;
  short nlink;
  uint size;
  uint addrs[NDIRECT + 1];
};
\end{lstlisting}

\subsection{Observations}
\begin{enumerate}[label=\arabic*.]
\item After \texttt{ln README r2}, both names had the same inode number \textbf{2}, the same size \textbf{2441 bytes}, and their link count became \textbf{2}. After \texttt{rm README}, \texttt{r2} still referred to inode 2 with the same size, while the link count decreased to \textbf{1}. Thus, the name changed, but the inode and data did not.
\item In \texttt{sys\_unlink}, the link count of the file inode is decremented by \texttt{ip->nlink--}. The file's blocks are not freed merely because one name is removed; they can be reclaimed only when the inode no longer has directory links and its final in-memory reference is released.
\item The xv6 on-disk inode contains the fields \texttt{type}, \texttt{major}, \texttt{minor}, \texttt{nlink}, \texttt{size}, and \texttt{addrs[NDIRECT + 1]}. Unlike the Linux experiment from Lab 5, the xv6 inode used here does not contain Unix permission bits, and this xv6 version has no symbolic-link type in \texttt{kernel/stat.h}.
\end{enumerate}

\section{Part 3. Measuring File-System Limits}

The constants controlling the basic file-system limits were read from \texttt{kernel/fs.h}:
\begin{lstlisting}
$ grep -n "define BSIZE\|define NDIRECT\|define NINDIRECT\|define DIRSIZ" kernel/fs.h
5:#define BSIZE   1024 // block size
26:#define NDIRECT   12
27:#define NINDIRECT (BSIZE / sizeof(uint))
54:#define DIRSIZ 14
\end{lstlisting}

A clean file-system image produced the following layout information:
\begin{lstlisting}
$ make clean
$ make fs.img 2>&1 | grep nmeta
nmeta 47 (boot, super, log blocks 31, inode blocks 13,
bitmap blocks 1) blocks 1953 total 2000
\end{lstlisting}

The name-length limit was measured in xv6:
\begin{lstlisting}
$ mkdir averyveryverylongname
$ ls
...
averyveryveryl 1 28 32
\end{lstlisting}

\subsection{Observations}
\begin{enumerate}[label=\arabic*.]
\item The requested name \texttt{averyveryverylongname} appears as \texttt{averyveryveryl}. The stored name is limited to \textbf{14 characters} by \texttt{DIRSIZ = 14}.
\item A file inode has 12 direct block addresses and one indirect block. With \texttt{BSIZE = 1024} bytes and a 4-byte \texttt{uint}, the indirect block can contain
\[
\mathrm{NINDIRECT} = \frac{1024}{4} = 256
\]
block addresses. Therefore the largest file is
\[
(12 + 256) \times 1024 = 274432 \text{ bytes} = 268 \text{ KiB}.
\]
\item The complete file-system image contains \textbf{2000 blocks}. The build reports \textbf{47 metadata blocks}: 1 boot block, 1 superblock, 31 log blocks, 13 inode blocks, and 1 bitmap block. The remaining \textbf{1953 blocks} are data blocks.
\end{enumerate}

\section{Part 4. The Recursive \texttt{find} Program}

The second user program, \texttt{user/find.c}, recursively walks a directory tree. It opens the requested directory, reads one \texttt{struct dirent} at a time, skips unused entries and \texttt{.}/\texttt{..}, builds each full path, compares the last component with the requested name, and recursively descends into subdirectories.

\begin{lstlisting}[language=C]
#include "kernel/types.h"
#include "kernel/stat.h"
#include "kernel/fs.h"
#include "user/user.h"
#include "kernel/fcntl.h"

void
find(char *path, char *target)
{
    char buf[512], *p;
    int fd;
    struct dirent de;
    struct stat st;

    if((fd = open(path, O_RDONLY)) < 0){
        fprintf(2, "find: cannot open %s\n", path);
        return;
    }

    if(fstat(fd, &st) < 0){
        fprintf(2, "find: cannot stat %s\n", path);
        close(fd);
        return;
    }

    if(st.type != T_DIR){
        close(fd);
        return;
    }

    if(strlen(path) + 1 + DIRSIZ + 1 > sizeof(buf)){
        fprintf(2, "find: path too long\n");
        close(fd);
        return;
    }

    strcpy(buf, path);
    p = buf + strlen(buf);
    if(p > buf && p[-1] != '/')
        *p++ = '/';

    while(read(fd, &de, sizeof(de)) == sizeof(de)){
        if(de.inum == 0)
            continue;

        char name[DIRSIZ + 1];
        memmove(name, de.name, DIRSIZ);
        name[DIRSIZ] = 0;

        if(strcmp(name, ".") == 0 || strcmp(name, "..") == 0)
            continue;

        memmove(p, de.name, DIRSIZ);
        p[DIRSIZ] = 0;

        if(stat(buf, &st) < 0)
            continue;

        if(strcmp(name, target) == 0)
            printf("%s\n", buf);

        if(st.type == T_DIR)
            find(buf, target);
    }

    close(fd);
}
\end{lstlisting}

The Makefile was updated with both new programs:
\begin{lstlisting}
$U/_mystat\
$U/_find\
\end{lstlisting}

A clean build successfully compiled \texttt{user/find.c} and produced \texttt{user/\_find}. The required recursive test was then performed:
\begin{lstlisting}
$ mkdir a
$ mkdir a/b
$ echo x > a/b/r2
$ echo y > a/r2
$ ln cat r2
$ find . r2
./a/b/r2
./a/r2
./r2
\end{lstlisting}

A second test searched for the directory named \texttt{b} below \texttt{a}:
\begin{lstlisting}
$ find a b
a/b
\end{lstlisting}

\subsection{Observations}
The output confirms that \texttt{find} traverses nested directories recursively and prints every matching path. The program does not recurse into \texttt{.} or \texttt{..}, preventing an infinite cycle. It also skips entries whose inode number is zero and closes each directory descriptor after use, which is important because xv6 provides only a small number of open file descriptors per process.

\section{Conclusion}
This laboratory work demonstrated how the xv6 file system implements familiar Unix concepts using a small and readable code base. Directory listings exposed the direct relationship between names, inode numbers, file types, and sizes. The \texttt{mystat} experiment showed that hard links are multiple directory names for the same inode and that unlinking one name only decreases the inode's link count.

The file-system constants revealed xv6's deliberate limits: 14-character directory names, 1024-byte blocks, 12 direct blocks, and one indirect block, giving a maximum file size of 268 KiB. The disk-image layout also showed the separation between metadata and data blocks. Finally, the recursive \texttt{find} program combined directory entries, inode metadata, system calls, path construction, and recursion into a practical file-system utility. Overall, the work connected xv6 source code directly to observable file-system behavior.

\end{document}
